%%%%%%%%%%%%%%%%%%%%%%%%%%%%%%%%%%%%%%%%%%%%%%%%%%%%%%%%%%%%%%%%%%%%%%
% Lớp 8 | Chương IV. Định lí Thalès | Tính chất đường phân giác của tam giác
% Dạng 8H4?17-2: Chứng minh tỉ lệ thức, bài toán tổng hợp
%%%%%%%%%%%%%%%%%%%%%%%%%%%%%%%%%%%%%%%%%%%%%%%%%%%%%%%%%%%%%%%%%%%%%%


% nguồn: De_AI_Toan8_cuoi_hoc_ki_i_1 (nhập 21/11/2026)
\begin{bt}%[8H4V17-2]%[ma:39a2ecdc]
	Cho $\triangle ABC$ có đường trung tuyến $AM$. Tia phân giác của góc $AMB$ cắt cạnh $AB$ tại $D$, tia phân giác của góc $AMC$ cắt cạnh $AC$ tại $E$.
a) Chứng minh $\dfrac{AD}{DB} = \dfrac{AE}{EC}$.
b) Chứng minh $DE \parallel BC$.
c) Gọi $I$ là giao điểm của $AM$ và $DE$. Chứng minh $I$ là trung điểm của $DE$.
	\loigiai{
		a) Xét $\triangle AMB$ có $MD$ là tia phân giác của góc $AMB$, theo tính chất đường phân giác ta có $\dfrac{AD}{DB} = \dfrac{MA}{MB}$. \\
Xét $\triangle AMC$ có $ME$ là tia phân giác của góc $AMC$, theo tính chất đường phân giác ta có $\dfrac{AE}{EC} = \dfrac{MA}{MC}$. \\
Vì $AM$ là đường trung tuyến của $\triangle ABC$ nên $MB = MC$, suy ra $\dfrac{MA}{MB} = \dfrac{MA}{MC}$, do đó $\dfrac{AD}{DB} = \dfrac{AE}{EC}$. \\
b) Xét $\triangle ABC$ có $\dfrac{AD}{DB} = \dfrac{AE}{EC}$ ($D \in AB$, $E \in AC$), theo định lí Thalès đảo ta suy ra $DE \parallel BC$. \\
c) Vì $DI \parallel BM$ (do $DE \parallel BC$) nên theo hệ quả định lí Thalès trong $\triangle ABM$ ta có $\dfrac{DI}{BM} = \dfrac{AI}{AM}$. \\
Tương tự, trong $\triangle ACM$ có $IE \parallel CM$ nên $\dfrac{IE}{CM} = \dfrac{AI}{AM}$. \\
Suy ra $\dfrac{DI}{BM} = \dfrac{IE}{CM}$. Vì $BM = CM$ nên $DI = IE$, tức $I$ là trung điểm của $DE$.
	}
\end{bt}

% nguồn: De_AI_Toan8_cuoi_hoc_ki_i_1 (nhập 21/11/2026)
\begin{bt}%[8H4V17-2]%[ma:3d0996b2]
	Cho $\triangle ABC$ có $AB = 12\text{ cm}$, $AC = 18\text{ cm}$, $BC = 20\text{ cm}$. Tia phân giác của góc $A$ cắt cạnh $BC$ tại $D$. Từ $D$ kẻ đường thẳng song song với $AB$ cắt $AC$ tại $E$, kẻ đường thẳng song song với $AC$ cắt $AB$ tại $F$.
a) Chứng minh tứ giác $AEDF$ là hình thoi.
b) Tính độ dài cạnh của hình thoi $AEDF$.
	\loigiai{
		a) Tứ giác $AEDF$ có $DE \parallel AF$ (vì $DE \parallel AB$) và $DF \parallel AE$ (vì $DF \parallel AC$) nên $AEDF$ là hình bình hành. \\
Hình bình hành $AEDF$ có đường chéo $AD$ là tia phân giác của góc $EAF$ nên $AEDF$ là hình thoi. \\
b) Đặt độ dài cạnh của hình thoi $AEDF$ là $x\text{ (cm)}$ ($x > 0$), khi đó $AE = DE = x$ và $CE = AC - AE = 18 - x$. \\
Xét $\triangle ABC$ có $DE \parallel AB$, theo hệ quả của định lí Thalès ta có: \\
$\dfrac{DE}{AB} = \dfrac{CE}{CA} \Rightarrow \dfrac{x}{12} = \dfrac{18 - x}{18} \Rightarrow 18x = 12(18 - x) \Rightarrow 18x = 216 - 12x \Rightarrow 30x = 216 \Rightarrow x = 7{,}2$. \\
Vậy độ dài cạnh của hình thoi $AEDF$ là $7{,}2\text{ cm}$.
	}
\end{bt}

% nguồn: De_AI_Toan8_cuoi_hoc_ki_i_1 (nhập 21/11/2026)
\begin{bt}%[8H4V17-2]%[ma:d481b170]
	Cho $\triangle ABC$ có ba đường phân giác trong là $AD$, $BE$, $CF$ ($D \in BC$, $E \in CA$, $F \in AB$).
a) Chứng minh rằng $\dfrac{DB}{DC} \cdot \dfrac{EC}{EA} \cdot \dfrac{FA}{FB} = 1$.
b) Áp dụng kết quả trên, tính độ dài các đoạn thẳng $FA$ và $FB$ biết $AB = 15\text{ cm}$, $\dfrac{DB}{DC} = \dfrac{3}{4}$ và $\dfrac{EC}{EA} = \dfrac{2}{3}$.
	\loigiai{
		a) Áp dụng tính chất đường phân giác trong $\triangle ABC$ cho ba đường phân giác $AD$, $BE$, $CF$, ta có: \\
$\dfrac{DB}{DC} = \dfrac{AB}{AC}$; $\dfrac{EC}{EA} = \dfrac{BC}{BA}$; $\dfrac{FA}{FB} = \dfrac{CA}{CB}$. \\
Nhân vế theo vế ba đẳng thức trên, ta được: \\
$\dfrac{DB}{DC} \cdot \dfrac{EC}{EA} \cdot \dfrac{FA}{FB} = \dfrac{AB}{AC} \cdot \dfrac{BC}{BA} \cdot \dfrac{CA}{CB} = 1$ (điều phải chứng minh). \\
b) Thay $\dfrac{DB}{DC} = \dfrac{3}{4}$ và $\dfrac{EC}{EA} = \dfrac{2}{3}$ vào hệ thức trên, ta có: \\
$\dfrac{3}{4} \cdot \dfrac{2}{3} \cdot \dfrac{FA}{FB} = 1 \Rightarrow \dfrac{1}{2} \cdot \dfrac{FA}{FB} = 1 \Rightarrow \dfrac{FA}{FB} = 2$. \\
Vì $FA + FB = AB = 15\text{ (cm)}$ và $FA = 2FB$, nên $2FB + FB = 15 \Rightarrow 3FB = 15 \Rightarrow FB = 5\text{ (cm)}$ và $FA = 10\text{ (cm)}$.
	}
\end{bt}

% nguồn: De_AI_Toan8_cuoi_hoc_ki_i_1 (nhập 21/11/2026)
\begin{bt}%[8H4V17-2]%[ma:84ae7cd3]
	Cho $\triangle ABC$, tia phân giác của góc $A$ cắt cạnh $BC$ tại $D$. Qua $D$ kẻ đường thẳng song song với $AB$ cắt cạnh $AC$ tại $E$.
a) Chứng minh $\triangle ADE$ cân tại $E$ và $\dfrac{AE}{EC} = \dfrac{AB}{AC}$.
b) Chứng minh $\dfrac{1}{DE} = \dfrac{1}{AB} + \dfrac{1}{AC}$.
	\loigiai{
		a) Vì $DE \parallel AB$ nên $\widehat{ADE} = \widehat{BAD}$ (hai góc so le trong). \\
Mặt khác $AD$ là tia phân giác của góc $A$ nên $\widehat{BAD} = \widehat{DAE}$. \\
Suy ra $\widehat{ADE} = \widehat{DAE}$, do đó $\triangle ADE$ cân tại $E$, suy ra $AE = DE$. \\
Xét $\triangle ABC$ có $DE \parallel AB$, theo định lí Thalès ta có $\dfrac{AE}{EC} = \dfrac{BD}{DC}$. \\
Lại có $AD$ là phân giác của góc $A$ nên $\dfrac{BD}{DC} = \dfrac{AB}{AC}$. Suy ra $\dfrac{AE}{EC} = \dfrac{AB}{AC}$. \\
b) Theo hệ quả của định lí Thalès trong $\triangle ABC$ có $DE \parallel AB$, ta có: \\
$\dfrac{DE}{AB} = \dfrac{CE}{AC} = \dfrac{AC - AE}{AC} = 1 - \dfrac{AE}{AC}$. \\
Vì $AE = DE$ nên $\dfrac{DE}{AB} = 1 - \dfrac{DE}{AC} \Rightarrow \dfrac{DE}{AB} + \dfrac{DE}{AC} = 1$. \\
Chia cả hai vế cho $DE > 0$, ta được $\dfrac{1}{AB} + \dfrac{1}{AC} = \dfrac{1}{DE}$ (điều phải chứng minh).
	}
\end{bt}

% nguồn: De_AI_Toan8_cuoi_hoc_ki_i_1 (nhập 21/11/2026)
\begin{bt}%[8H4V17-2]%[ma:01b5c8ec]
	Cho hình bình hành $ABCD$ có $AB < AD$ và hai đường chéo $AC$, $BD$ cắt nhau tại $O$. Tia phân giác của góc $BAD$ cắt đường chéo $BD$ tại $M$, tia phân giác của góc $ABC$ cắt đường chéo $AC$ tại $N$.
a) Chứng minh $\dfrac{MB}{MD} = \dfrac{NA}{NC}$.
b) Chứng minh $MN \parallel CD$.
	\loigiai{
		a) Xét $\triangle ABD$ có $AM$ là tia phân giác của góc $BAD$, theo tính chất đường phân giác ta có $\dfrac{MB}{MD} = \dfrac{AB}{AD}$. \\
Xét $\triangle ABC$ có $BN$ là tia phân giác của góc $ABC$, theo tính chất đường phân giác ta có $\dfrac{NA}{NC} = \dfrac{BA}{BC}$. \\
Vì $ABCD$ là hình bình hành nên $AD = BC$, suy ra $\dfrac{MB}{MD} = \dfrac{NA}{NC}$. \\
b) Từ $\dfrac{MB}{MD} = \dfrac{NA}{NC}$ suy ra $\dfrac{MB}{MB + MD} = \dfrac{NA}{NA + NC}$, hay $\dfrac{MB}{BD} = \dfrac{NA}{AC}$. \\
Vì $AB < AD$ nên $MB < MD$, do đó $MB < \dfrac{1}{2}BD = OB$, tức điểm $M$ nằm giữa $O$ và $B$. Tương tự, điểm $N$ nằm giữa $O$ và $A$. \\
Do $OB = \dfrac{1}{2}BD$ và $OA = \dfrac{1}{2}AC$ nên $\dfrac{MB}{OB} = \dfrac{2MB}{BD} = \dfrac{2NA}{AC} = \dfrac{NA}{OA}$. \\
Suy ra $\dfrac{OB - MB}{OB} = \dfrac{OA - NA}{OA}$, hay $\dfrac{OM}{OB} = \dfrac{ON}{OA}$. \\
Xét $\triangle OAB$ có $\dfrac{OM}{OB} = \dfrac{ON}{OA}$ ($M \in OB$, $N \in OA$), theo định lí Thalès đảo ta có $MN \parallel AB$. \\
Vì $AB \parallel CD$ nên suy ra $MN \parallel CD$.
	}
\end{bt}
